\chapter{Functional Operations}\label{ch:functional}

Functions in Dylan are objects: they can be passed as arguments, stored
in variables and collections, and returned as results
(chapter~\ref{ch:functions}). The functions in this chapter make new
functions out of existing ones. They come from chapter~17 of the 1992
manual (\dylanbook{chapter 17}), with \code{identity} from chapter~18
and \code{apply} from chapter~8. Combined with the collection functions
of chapter~\ref{ch:collections}, they let many loops be written as a
single expression.

Each function here answers a new function that accepts any number of
arguments (as many as the functions it was made from need). In both
implementations they are ordinary closures; the examples show their
results as comments, and all of them give the same results in the
compiler to C and the Common Lisp implementation.

\section{compose}\label{sec:functional-compose}

\index{compose}
\code{(compose \meta{f} \meta{g} ...)} answers a function that calls the
last function on its arguments, then each function before it on the
previous result, from right to left. With one argument, \code{compose}
answers that argument. Thus \code{cadr} could be written
\code{(compose car cdr)}:

\begin{dylan}
((compose car cdr) '(1 2 3))              ; => 2
((compose (curry * 2) (curry + 1)) 5)     ; => 12     (5 + 1) * 2
(map (compose - abs) '(-1 2 -3))          ; => (-1 -2 -3)
\end{dylan}

\noindent Only the first value of each inner call is passed on:
\code{((compose list floor) 7/2)} is \code{(3)}, though \code{floor}
answers two values.

This is the manual's example: the distance of a point from the origin,
built from three small functions.

\begin{dylan}
(define-method sum ((numbers <sequence>))
  (reduce1 + numbers))
(define-method square ((x <number>)) (* x x))
(define-method square-all ((coords <sequence>))
  (map square coords))

(define distance (compose sqrt sum square-all))
(distance '(3 4 5))      ; => 7.0710678118654755
\end{dylan}

\section{complement, disjoin and conjoin}\label{sec:functional-logic}

\index{complement}
\code{(complement \meta{predicate})} answers a function that answers
\code{\#t} where the predicate answers \code{\#f}, and \code{\#f}
otherwise:

\begin{dylan}
(map (complement even?) '(1 2 3))          ; => (#t #f #t)
((complement id?) 'a 'b)                   ; => #t
((complement member?) 5 '(1 2))            ; => #t
(define not-empty? (complement empty?))
(choose not-empty? '("a" "" "bc"))         ; => ("a" "bc")
\end{dylan}

\index{disjoin}
\code{(disjoin \meta{p} \meta{q} ...)} answers the \term{disjunction}
of its predicates: a function that calls them in order on its arguments
and answers the first true value any of them returns, without calling
the rest, or \code{\#f} if none does. It is like an \code{or} of calls.

\begin{dylan}
(define small-or-even? (disjoin (curry > 3) even?))
(choose small-or-even? '(1 2 3 4 5 6 7))           ; => (1 2 4 6)
((disjoin (method (x) #f) (method (x) (* x 10))) 4) ; => 40
(choose (disjoin zero? (curry < 10)) '(0 5 11 3 20)) ; => (0 11 20)
\end{dylan}

\index{conjoin}
\code{(conjoin \meta{p} \meta{q} ...)} answers the \term{conjunction}:
a function that calls the predicates in order and answers \code{\#f} as
soon as one does, or else the value of the last. It is like an
\code{and} of calls.

\begin{dylan}
(define positive-integer? (conjoin (rcurry instance? <integer>) positive?))
(map positive-integer? '(5 -5 5.0))                    ; => (#t #f #f)
((conjoin (rcurry instance? <integer>) (method (x) (* x 2))) 4)   ; => 8
(find-key '(3 8 12 5) (conjoin even? (curry < 10)))    ; => 2
\end{dylan}

\section{curry and rcurry}\label{sec:functional-curry}

\index{curry}\index{rcurry}
\code{(curry \meta{function} \meta{arg} ...)} answers a function that
calls \meta{function} with the given arguments followed by its own.
\code{rcurry}, ``right curry'', puts the given arguments after its own
instead:

\begin{dylan}
((curry list 1 2) 3 4)             ; => (1 2 3 4)
((rcurry list 1 2) 3 4)            ; => (3 4 1 2)
(map (curry + 1) '(3 4 5))         ; => (4 5 6)
(map (rcurry - 1) '(5 6 7))        ; => (4 5 6)
\end{dylan}

\noindent So \code{(curry id? 'x)} tests for identity with the symbol
\code{x}; \code{(curry > 6)} is true of numbers less than 6, and
\code{(rcurry > 6)} of numbers greater than 6:

\begin{dylan}
((curry id? 'x) 'x)                ; => #t
((curry > 6) 3)                    ; => #t
((rcurry > 6) 3)                   ; => #f
((curry concatenate "set-") "size" "!")    ; => "set-size!"
(define yuppify (rcurry concatenate ", ayup"))
(yuppify "I'm from New Hampsha")   ; => "I'm from New Hampsha, ayup"
\end{dylan}

\noindent \code{rcurry} is the natural way to fix the second argument of
a two-argument function, as in \code{(rcurry instance? <integer>)}, a
predicate for integers, or \code{(rcurry element 0)}, which takes the
first element of a sequence:

\begin{dylan}
(map (rcurry element 0) '("abc" "xyz"))    ; => (#\a #\x)
\end{dylan}

\section{always and identity}\label{sec:functional-always}

\index{always}
\code{(always \meta{object})} answers a function that ignores any
arguments and answers the object:

\begin{dylan}
((always 1) 'x 'y 'z)       ; => 1
((always #t) #f #f)         ; => #t
((always 'k))               ; => k
(map (always 0) '(a b c))   ; => (0 0 0)
\end{dylan}

\index{identity}
\code{(identity \meta{object})} answers its argument. It is useful
where a function is required and nothing should be done, such as the
start of a chain of compositions:

\begin{dylan}
(identity 5)                ; => 5
(map identity '(1 2))       ; => (1 2)
(define add1 (curry + 1))
((reduce (method (f g) (compose f g)) identity (list add1 add1 add1)) 0)
; => 3
\end{dylan}

\section{apply}\label{sec:functional-apply}

\index{apply}
\code{(apply \meta{function} \meta{arg} ... \meta{sequence})} calls the
function with the given arguments followed by the elements of the
sequence, which may be any sequence (\dylanbook{chapter 8}):

\begin{dylan}
(apply + 1 2 '(3 4))                    ; => 10
(apply max '(3 1 4))                    ; => 4
(apply list '())                        ; => ()
(apply + #(1 2 3))                      ; => 6
(apply + 1 (range from: 1 through: 3))  ; => 7
(apply (curry + 100) '(1 2))            ; => 103
(apply concatenate '("a" "b" "c"))      ; => "abc"
\end{dylan}

\noindent Since \code{map} takes any number of collections, applying it
to a list of lists transposes them:

\begin{dylan}
(apply map list '((1 2 3) (a b c)))     ; => ((1 a) (2 b) (3 c))
\end{dylan}

\section{Idioms}\label{sec:functional-idioms}

\index{idioms}
These operations are at their best as the arguments of \code{map},
\code{choose}, \code{reduce}, \code{any?}, \code{every?} and
\code{find-key}, where they replace small anonymous methods.

\paragraph{Filtering and transforming.} Choose the odd numbers from 1 to
10, double them, and add them up:

\begin{dylan}
(reduce + 0 (map (curry * 2) (choose odd? (range from: 1 through: 10))))
; => 50
\end{dylan}

\paragraph{Negating a test.} Keep the lists that do not contain 2.
Because \code{member?} takes the value first, \code{curry} fixes it:

\begin{dylan}
(choose (complement (curry member? 2)) '((1 2) (3 4) (2 5)))
; => ((3 4))
\end{dylan}

\paragraph{Testing membership in a fixed set.} \code{rcurry} fixes the
collection instead. On a list \code{member?} answers the tail where the
element was found, which \code{any?} passes on as its true value:

\begin{dylan}
(any? (rcurry member? '(a b)) '(x y b))     ; => (b)
\end{dylan}

\paragraph{Counting.} Count the odd numbers greater than 3:

\begin{dylan}
(define count-if (method (pred seq) (size (choose pred seq))))
(count-if (conjoin odd? (curry < 3)) '(1 2 3 4 5 6 7))   ; => 2
\end{dylan}

\paragraph{Comparing by a key.} Sort strings without regard to case,
and find the length of the longest:

\begin{dylan}
(sort '("pear" "Apple" "fig")
      test: (method (a b) (< (as-lowercase a) (as-lowercase b))))
; => ("Apple" "fig" "pear")
(reduce1 max (map size '("pear" "banana" "fig")))    ; => 6
\end{dylan}

\paragraph{Checking several classes.} \code{curry} fixes the object and
leaves the class to vary:

\begin{dylan}
(every? (curry instance? 3) (list <integer> <rational> <number>))   ; => #t
\end{dylan}
